\documentclass[11pt]{article}
\usepackage{amsmath,amssymb,amsthm}

\title{The Bollob\'as--Leader Conjecture on Few Subsequence Sums in $\mathbb Z_n\times \mathbb Z_n$: A Partial Result}
\author{Autonomous Mathematical Research Project}
\date{\today}

\newtheorem{theorem}{Theorem}
\newtheorem{lemma}{Lemma}
\newtheorem{conjecture}{Conjecture}

\begin{document}
\maketitle

\begin{abstract}
We study the Bollob\'as--Leader conjecture on the minimum number of distinct
nonempty subsequence sums of a zero-free sequence in $(\mathbb Z/n\mathbb Z)^2$.
We give a complete proof of the case $t=0$: every zero-free sequence of
length $n-1$ has at least $n-1$ distinct nonempty subsequence sums.  We also
report exhaustive computational verification of the conjectured bound for
$n=3$ (all $t$), $n=4$ (all $t$), and $n=5$ ($t=0,1,2$).  The general
conjecture remains open; the obstruction to a direct induction is identified.
\end{abstract}

\section{Problem statement}

Let $G_n=(\mathbb Z/n\mathbb Z)^2$.  A sequence $(a_1,\dots,a_\ell)\in G_n^\ell$ is
\emph{zero-free} if no nonempty subsequence sums to $0\in G_n$.  A
\emph{subsequence sum} is any sum $a_{i_1}+\cdots+a_{i_r}$ with
$1\le i_1<\cdots<i_r\le \ell$.

\begin{conjecture}[Bollob\'as--Leader]
For $0\le t\le n-1$, every zero-free sequence in $G_n$ of length $n-1+t$
has at least $n(t+1)-1$ distinct nonempty subsequence sums.
\end{conjecture}

The extremal construction is $n-1$ copies of $(1,0)$ together with $t$
copies of $(0,1)$, which is zero-free and attains $n(t+1)-1$ distinct
subsequence sums.

\section{Existing results}

\begin{itemize}
\item Olsen's theorem: every sequence of length at least $2n-1$ in
$G_n$ contains a nonempty zero-sum subsequence; hence zero-free
sequences have length at most $2n-2$.
\item One-dimensional analogue: in $\mathbb Z_n$, the sequence of $t$
copies of $1$ minimises the number of distinct nonempty subsequence
sums among zero-free sequences of length $t$.
\end{itemize}

\section{Result proved here}

\begin{theorem}
If $\mathbf a=(a_1,\dots,a_{n-1})$ is a zero-free sequence in $G_n$, then
$\mathbf a$ has at least $n-1$ distinct nonempty subsequence sums.
\end{theorem}

\begin{proof}
Define $S_0=0$ and $S_i=a_1+\cdots+a_i$ for $1\le i\le n-1$.  If
$S_i=S_j$ with $0\le i<j\le n-1$, then
$a_{i+1}+\cdots+a_j=0$, a nonempty subsequence summing to $0$,
contradicting zero-freeness.  Thus $S_0,\dots,S_{n-1}$ are $n$ distinct
elements.  Removing $S_0=0$, the remaining $n-1$ values are distinct
nonzero elements of $\Sigma^\circ(\mathbf a)$.  Hence
$|\Sigma^\circ(\mathbf a)|\ge n-1$.
\end{proof}

This is the case $t=0$ of the conjecture.

\section{Computational verification}

The file \texttt{src/min\_sums.py} exhaustively checks the conjectured
minimum for $n=3$ (all $t$), $n=4$ (all $t$), and $n=5$ ($t=0,1,2$);
see \texttt{experiments/computation\_results.md}.  In every tested
case the minimum agrees with $n(t+1)-1$.

\section{Obstruction to a full proof}

A natural induction would say: appending one element to a zero-free
sequence increases the number of distinct subsequence sums by at least $n$.
Zero-freeness rules out $-x$ being a previous sum, but it does not rule
out other collisions between the old set and its translate by $x$.  A
collision $y=x+z$ with $y,z$ previous sums need not yield a zero-sum
subsequence, because the two underlying subsequences may overlap
nontrivially and common terms need not cancel in subset-sum notation.
Controlling these collisions is the remaining obstacle.

\section{Conclusion}

The Bollob\'as--Leader conjecture is proved here only for $t=0$.  The
general case remains open.  The supplied code verifies the predicted
minima for small $n$.

\end{document}
